\section{Conclusion}
\label{sec:conclusion}

The marginal value of surveillance case volume is governed by the interplay between the inferential target and intrinsic error floors. When monitoring current transmission potential, returns to sample volume do not saturate, and surveillance capacity should be determined by statistical detection power. When estimating cross-period transmission trends, expanding single-period case counts yields diminishing returns; once sample volume reaches the variance balance scale $m_\times$, additional resources should be redirected toward accumulating observation periods, expanding spatial pooling, standardizing reporting protocols, and regularly updating baseline benchmark models. Empirical analysis of US influenza hospitalisations confirms an error scaling pattern compatible with this theoretical decomposition, in which unmodelled error components account for approximately half of early-epidemic loss and cannot be explained by reporting completeness proxies, though current observations cannot separate a strict non-zero asymptote from slow decay. Resolving the scale of environmental variance requires accumulating observation periods across multiple epidemic seasons. Surveillance systems must therefore move beyond single-period sample volume as the sole metric of monitoring adequacy.
